\section{相关工作}

本文压在四支文献的交汇处；每支既给口径，也给出本篇要绕开的那一步。

\textbf{复杂网络的骨架一支。}
随机图给了"平均度数为常"的基线，而 \citet{barabasi1999} 指出真实网络的度数分布
是无标度的——少数枢纽扛起绝大部分连接。本篇正落在这一支里，
但只取它的一半结论：枢纽确实存在（最大度数 113，单人占全部度数和的 2.5\%），
本篇不去论证"因此是无标度"，而去检验这个标签在数据上是否站得住（第 5.2 节）。
\citet{watts1998} 的小世界与 \citet{albert2000} 的攻击容忍度进一步说明：
一张看起来很密的网，可以只靠极少数节点连着，抽掉它们就散。
本篇第 5.3 节报的"18 个连通块、最大块占带边顶点的 94.9\%"，就是这条警告的直接读数。

\textbf{幂律拟合的边界一支，本篇主要借这一支。}
\citet{clauset2009} 给出了幂律拟合的完整判据：尾指数必须在大样本、多切点上稳定，
$x_{\min}$ 的选取要能自证，拟合优度要用连续检验而非分箱比较。
本篇第 5.2 节完全按这套口径去做，得到的是一条负结果：
尾指数随 xmin 漂移 0.971、尾样本只有 142 条，故不贴幂律标签。
把"复杂网络的常见形状"直接写进历史叙事研究的做法，
在没有走这一道检验之前，本篇一律不采用。
\textbf{网络方法进入历史学的一支。}
\citet{lemercier2015} 论证历史学者如何用网络方法，并点明首要陷阱：
从历史文本抽出的关系是叙述者建构的关系，不是被研究时代里发生过的关系。
本篇因此不复现任何"明代社会真实结构"的论断，只报这部叙事自己长成的形状。
\textbf{历史动力学一支，本篇借它作对照假设而不作证据。}
\citet{turchin2003} 主张历史宏观过程可用人口—财政的周期模型解释，
\citet{diamond2005} 则从多个社会比较"选择失败还是存活"。
两支都给了一个可对照的直觉：崩塌是结构性位移，不是单次事件。
本篇能验的只有其中一个可数项——结构位移在时间上是否真的发生了（第 5.5 节的类别迁移），
而且验的是文本抽取出来的事件类别，不是人口或财政数据。
方法层面本量级（1231 个顶点、2324 条边）的网络综述口径取自
\citet{newman2005} 与 \citet{wasserman1994}；本篇的基尼、连通块、
秩相关一律按这两支的定义算。

\textbf{本文的位置。}
四支里没有一支把"王朝兴衰"这个对象上的四个量（集中、幂律边界、连通、时间迁移）
同时点数出来：网络那两支给了形状与判据，历史学那两支给了对照假设与陷阱，
中间缺一层——用一个具体语料把这四个量一次算齐，并显式负报其中不成立者。
本篇用同一份图谱一次补齐，且把"不主张幂律"与"样本量上限"写进正文而非脚注。
